% teaser % lin
% 3D 打印视频
% 渲染视频
\section{Introduction}

Part segmentation provides explicit part-level structures of 3D assets, serving as a core primitive for 3D content creation pipelines and offering fundamental 3D perception capabilities for spatial intelligence.
%
It enables a wide range of downstream applications, including part-level editing, animation rigging, and industrial uses such as 3D printing.
%
However, existing methods often fall short in segmentation quality, producing erroneous regions and imprecise boundaries that limit their practical usability.


To this end, one line of work attempts to transfer the comprehensive 2D segmentation priors to 3D via 2D-to-3D lifting.
%
Methods such as SAMPart3D~\cite{sampart3d} optimize 3D segmentation via 2D-to-3D distillation, 
but incur substantial computational and time overhead, and often yield blurry boundaries.
%
In parallel, another set of methods~\cite{segment3d,pointsam,geosam2} applies SAM~\cite{sam,sam2,sam3} to obtain 2D masks of multi-view projected images, which are then back-projected and fused into 3D masks.
%
However, these multi-view pipelines incur substantial runtime overhead, are sensitive to view coverage, and the back-projection and fusion step often introduces cross-view inconsistencies and imprecise boundaries.

% su du kuai
Recently, another line of work~\cite{p3sam,partsam} moves toward native 3D part segmentation so as to remedy the inherent shortcomings of the aforementioned methods that leverage 2D segmentation priors.
%
These methods predict segmentation parts directly in the native 3D space, explicitly enforce semantic and structural consistency, and are more efficient at inference.
%
However, it is a typical requirement to collect large-scale training datasets with curated 3D part annotations, 
where fine-grained annotations are costly and inconsistent across sources in granularity, hierarchy, and boundary definitions.


%
% 总结，2D的问题，mismatch，3D的问题，然后本质promising是纯三维，41行，我们发现能够物体结构，外观这里，分割，prior
% Overall, existing approaches face a fundamental trade-off between annotation and training cost and the structural artifacts caused by lifting 2D cues into 3D.
% %
% We attribute these artifacts to the inherent mismatch between 2D priors and native 3D structure,
% %
% motivating us to instead exploit a native 3D structural prior in place of lifted 2D cues.
In summary, the first line of methods suffers from a mismatch between 2D priors and 3D structure, while the second relies on costly training from scratch.
%
Therefore, a more promising approach is to leverage a prior model that encodes both 3D structure and texture to perform segmentation.
%
In particular, 3D generative models trained on large-scale unannotated 3D textured assets internalize rich part-level structure and texture patterns, providing a strong 3D prior over geometry and appearance. 
%
Such priors encourage part segmentation with sharper boundaries, while reducing reliance on dense part annotations and extensive task-specific training.
%
This motivates us to ask:
\textit{How can 3D generative priors be effectively transferred to part-level 3D segmentation to improve quality and data efficiency?}


Motivated by this perspective, we propose \textit{SegviGen}, a generative framework for 3D part segmentation that leverages the rich 3D structural and textural knowledge encoded in large-scale 3D generative models.
%
Specifically, we formulate part segmentation as a colorization task that fully exploits the capacity of 3D generative models.
%
\textit{SegviGen} encodes the input 3D asset into a latent representation and uses it, together with the task embedding and query points, to condition the denoising process.
%
The model is trained to predict part-indicative colors, along with reconstructing the underlying geometry.
%
This formulation naturally accommodates additional conditioning signals, enabling \textit{SegviGen} to flexibly support interactive part segmentation, full segmentation, and 2D segmentation map–guided full segmentation under a unified architecture. 
%
Notably, while the first two tasks are common settings, 2D segmentation map–guided full segmentation is uniquely enabled by \textit{SegviGen}, supporting arbitrary part granularity and more precise segmentation that is critical for industrial applications.
% 写清楚哪几个setting


Qualitative and quantitative results show that \textit{SegviGen} consistently surpasses the prior state of the art, P3-SAM~\cite{p3sam}, while using only 0.32\% of the labeled training data. % 写清楚啥方法。
%
On interactive part segmentation, it achieves the best performance across all metrics on PartObjaverse-Tiny~\cite{partobjaverse_tiny} and PartNeXT~\cite{partnext}, with a 40\% gain in IoU@1, an important metric that reflects the model’s single-click accuracy. 
% 体现这个指标的挑战性 
%
On full segmentation without guidance, \textit{SegviGen} outperforms the best baseline by 15\% in overall IoU, averaged across datasets.
%
Our main contributions are summarized as follows:
\begin{itemize}[leftmargin=*, topsep=2pt, itemsep=1pt, parsep=0pt, partopsep=0pt]
    \item We propose \textit{SegviGen}, a unified multi-task framework for 3D part segmentation that effectively exploits the structural and textural priors encoded in pretrained 3D generative models, enabling accurate and efficient segmentation.
    \item We reformulate 3D segmentation as part-wise colorization, where \textit{SegviGen} predicts the colors of actiave voxel as part labels in a single generative process.
    \item Extensive experiments show that \textit{SegviGen} outperforms the prior state of the art by 40\% on interactive part segmentation and 15\% on full segmentation, using only 0.32\% of the labeled training data, highlighting the effectiveness of transferring 3D generative priors to part segmentation.
\end{itemize}

